% QUESTAO
Após determinar a base e a altura do triângulo abaixo, a sua área, em metros quadrados, fica:

\parbox{0.4\linewidth}{

% ALTERNATIVAS
$9\sqrt{5}$ m$^2$
$18\sqrt{5}$ m$^2$
$4\sqrt{3}$ m$^2$
$9\sqrt{2}$ m$^2$
$18\sqrt{3}$ m$^2$

}\parbox{0.59\linewidth}{
\begin{center}
\vspace{-6mm}
\begin{tikzpicture}[ >=latex, scale=0.4 ]

\coordinate (A) at (0,0);
\coordinate (B) at (9,0);
\coordinate (C) at (4,2*sqrt(5));
\coordinate (D) at (4,0);

\draw[dashed] (C) -- (D);
\pgfmathsetmacro{\raio}{1};
\draw (D) rectangle +(0.6*\raio,0.6*\raio);
\fill ($(D)+(0.3*\raio,0.3*\raio)$) circle (0.5*\raio mm);

\draw[thick] (A) -- (B);
\draw[thick,nodelabel={-12pt}{$3\sqrt{5}$}] (B) -- (C);
\draw[thick,nodelabel={-8pt}{$6$}] (C) -- (A);
\draw[nodelabel={8pt}{$4$}] (A) -- (D);

\end{tikzpicture}
\vspace{-9mm}
\end{center}
}


